\section{Riemannian \methodname{}}\label{sec:riemannian}

This section provides additional details on the extension of TDS to Riemannian diffusion models introduced in \Cref{sec:method}.
We first introduce the tangent normal distribution.
We then provide with background on Riemannian diffusion models, which we parameterize with the tangent normal.
Then we describe the extension of TDS to these models.
Finally we show how \Cref{alg:TDS} modifies to this setting.

\paragraph{The tangent normal distribution.}
Just as the Gaussian is the parametric family underlying Euclidean diffusion models in \Cref{eq:unconditional-diffusion-transition}, 
the tangent normal (see e.g.\ \citep{chirikjian2014gaussian,deriemannian})
underlies generation in Riemannian diffusion models so we review it here.

We take the tangent normal to be the distribution is implied by a two step procedure.
Given a variable $x$ in the manifold, the first step is to sample a variable $\bar y$ in $\mathcal{T}_{x},$ the tangent space at $x;$
if $\{h_1, \dots, h_D\}$ is an orthonormal basis of $\mathcal{T}_{x}$
one may generate
\begin{align}
    \bar y = \mu + \sum_{d=1}^D \sigma \epsilon_d \cdot h_d,
\end{align}
with $\epsilon_d\overset{i.i.d.}{\sim}\mathcal{N}(0, 1),$
and $\sigma^2>0$ a variance parameter. 
The second step is to project $\bar y$ back onto the manifold to obtain
$y = \exp_{x}\{ \bar y\}$
where $\exp_{x}\{ \cdot \}$ denotes the exponential map at $x.$ The resulting distribution of $y$ is denoted as $\mathcal{TN}_{x}(\mu, \sigma^2)$. 
By construction, $\mathcal{T}_x$ is a Euclidean space with its origin at $x,$
so when $\|\mu\|_2=0, \mathcal{TN}_{x}(\mu, \sigma^2)$ is centered on $x.$
And since the geometry of a Riemannian manifold is locally Euclidean, when $\sigma^2$ is small
the exponential map is close to the identity map and the tangent normal is essentially a narrow Gaussian distribution in the manifold at $x.$
Finally, we use $\mathcal{TN}_{x}(y;\mu, \sigma^2)$ to denote the density of the tangent normal evaluated at $y$.

Because this procedure involves a change of variables from $\bar y$ to $y$ (via the exponential map),
to compute the tangent normal density one computes 
\begin{equation}\label{eq:TN_density}
    \mathcal{TN}_{x}(y;\mu, \sigma^2) = \mathcal{N}(\exp^{-1}_x\{y\}; \mu,  \sigma^2) \left| \frac{\partial}{\partial y} \exp_x^{-1}\{y\}\right|
\end{equation}
where $\exp_x^{-1}\{y\}$ is the inverse of the exponential map (from the manifold into $\mathcal{T}_x$), 
and $\left| \frac{\partial}{\partial y} \exp_x^{-1}\{y\}\right|$ is the determinant of the Jacobian of the exponential map; when the manifold lives in a higher dimensional subset of $R,$ we take 
$\left| \frac{\partial}{\partial y} \exp_x^{-1}\{y\}\right|$ to be the product of the positive singular values of the Jacobian.

\paragraph{Riemannian diffusion models and the tangent normal distribution.}
Riemannian diffusion models proceeds through a geodesic random walk \citep{deriemannian}.
At each step $t,$
one first samples a variable $\dt{\bar x}{t}$ in $\mathcal{T}_{\xtpo};$
if $\{h_1, \dots, h_D\}$ is an orthonormal basis of $\mathcal{T}_{\xtpo}$
one may generate
\begin{align}
    \dt{\bar x}{t}  =\sigma_{t+1}^2 s_\theta(\dt{x}{t+1})+ \sum_{d=1}^D \sigma_{t+1} \epsilon_d \cdot h_d,
\end{align}
with $\epsilon_d\overset{i.i.d.}{\sim}\mathcal{N}(0, 1).$
One then projects $\dt{\bar x}{t}$ back onto the manifold to obtain
$\xt = \exp_{\xtpo}\{ \dt{\bar x}{t}\}$
where $\exp_{x}\{ \cdot \}$ denotes the exponential map at $x.$

This is equivalent to sampling $\xt$ from a tangent normal distribution as 
\begin{equation}\label{eqn:reverse_transition_unconditional_riemannian}
    \dt{x}{t} \sim p(\dt{ x}{t} \mid \xtpo) = \mathcal{TN}_{\xtpo}\left(\dt{ x}{t};\sigma_{t+1}^2 s_\theta(\dt{x}{t+1}),
\sigma_{t+1}^2 \right).
\end{equation}

\paragraph{TDS for Riemannian diffusion models.}
To extend TDS, appropriate analogues of the twisted proposals and weights are all that is needed.
For this extension we require that the diffusion model is also associated with a manifold-valued denoising estimate $\denoiser$ as will be the case when, for example, $\scoremodel(\xt, t) := \nabla_{\xt} \log q(\xt \mid \xzero=\denoiser)$ for $\denoiser = \denoiser(\xt, t).$
In contrast to the Euclidean case, a relationship between a denoising estimate and a computationally tractable score approximation may not always exist for arbitrary Riemannian manifolds;
however for Lie groups when the the forward diffusion is the Brownian motion, tractable score approximations do exist \citep[Proposition 3.2]{yim2023se}.

For the case of positive and differentiable $\plikelihood{y}{\xzero},$ we again choose twisting functions $\twist(y\mid \xt):= \plikelihood{y}{\denoiser(\xt)}.$

Next are the inpainting and inpainting with degrees of freedom cases.
Here, assume that $\xzero$ lives on a multidimensional manifold (e.g.\ $SE(3)^N$)
and the unmasked observation $y=\xzero_\mask$ with $\mask \subset \{1, \dots, N\}$ on a lower-dimensional sub-manifold (e.g.\ $SE(3)^{|\mask|},$ with $|\mask|<N$).
In this case, twisting functions are constructed exactly as in \Cref{subsec:conditioning-problems}, 
except with the normal density in \Cref{eqn:inpaint_twist} replaced with a Tangent normal as
\begin{align}
\twist(y\mid \xt) {=} \mathcal{TN}_{\denoiser(\xt)_\mask}( y ; 0, \bar \sigma_t^2).
\end{align}

For all cases, we propose the twisted proposal as
\begin{equation}\label{eqn:likelihood_approx_riemannian_supp}
    \pTwist(\dt{ x}{t} \mid \xtpo, y) {=} \mathcal{TN}_{\xtpo}\left(\dt{ x}{t};\sigma_{t+1}^2 s_\theta(\dt{x}{t+1}, y), \tilde\sigma_{t+1}^2 \right)
\end{equation}
%where $\bar y := \exp_{\denoiser(\dt{x}{t})_\umask}^{-1}(y)\in \mathcal{T}_{\denoiser(\dt{x}{t})_\umask}$ 
%is $\dt{x}{0}$ mapped into $\mathcal{T}_{\denoiser(\dt{x}{t})}$ by the inverse of the exponential map at $\denoiser(\dt{x}{t})_\umask,$
where as in the Euclidean case  $s_\theta(\dt{x}{t}, y)=s_\theta(\dt{x}{t}) + \nabla_{\dt{x}{t}} \log \tilde p(y \mid \dt{x}{t})$.


Weights at intermediate steps are computed as in the Euclidean case (\Cref{eq:diffusion-twisted-weigthing-functions}): 
%However, 
%since the tangent normal includes a projection step back to the manifold, its density involves not only usual Gaussian density but also the Jacobian of the exponential map to account for the change of variables.
%For example, in the case that the ambient dimension of the manifold and the tangent space are the same dimension,
%$$
%\mathcal{TN}_{x}(x^\prime;\mu, \sigma^2)= \mathcal{N}(\exp^{-1}_x\{x^\prime\}; \mu, \sigma^2)\left|\frac{\partial}{\partial x} \exp_x^{-1}\{x^\prime\}\right|,
%$$
%where $\exp^{-1}_x\{x^\prime\}$ is the inverse of the exponential map at $x$ from the manifold into $\mathcal{T}_x,$
%and 
%$\left|\frac{\partial}{\partial x} \exp_x^{-1}\{x^\prime\}\right|$ is the determinant of its Jacobian.
%The weights are then computed as
\begin{align}\label{eq:diffusion-twisted-weigthing-functions-riemannian}
    \weight_t (\xt , \xtpo ):&= \frac{\pmodel (\xt \mid \xtpo) \twist(y \mid \xt)}{\twist(y \mid \xtpo) \pTwist(\xt \mid \xtpo, y)} \\
    &= \frac{\mathcal{TN}_{\xtpo}(\xt; \sigma_{t+1}^2\scoremodel(\xtpo), \sigma_{t+1}^2) \twist(y \mid \xt)
    }{\twist(y \mid \xtpo) \mathcal{TN}_{\xtpo}(\xt; \sigma_{t+1}^2\scoremodel(\xtpo, y), \tilde\sigma_{t+1}^2)}.
\end{align}
While the proposal and target contribute identical Jacobian determinant terms that cancel out, they remain in the twisting functions.


\paragraph{Adapting the TDS algorithm to the Riemannian setting.}
To translate the TDS algorithm to the Riemannian setting we require only two changes.

The first is on \Cref{alg:TDS} 
\Cref{alg_line:conditional_score}.
Here we assume that the unconditional score is computed through the transition density function: 
\begin{align}
\scoremodel(\xtpo) := \nabla_{\xtpo} \log q_{t+1{|}0}(\xtpo\mid \denoiser) \quad \mathrm{\ for\ } \denoiser = \denoiser(\xtpo).
\end{align}
Note that the gradient above ignores the dependence of $\denoiser$ on $\xtpo$. 

%\BLT{Is it clear that the gradient above ignores the dependence of $\denoiser$ on $\xt$?}
The conditional score approximation in \Cref{alg_line:conditional_score} is then replaced with 
$$
\tilde s_k \leftarrow \pTwist (\cdot \mid \xtpo_k, y):= \nabla_{\xtpo_k} \log q_{t+1{|} 0}(\xtpo_k\mid \denoiser) + \nabla_{\xtpo_k} \log \tilde p_k^{t+1} \quad \mathrm{ \ for\  } \denoiser = \denoiser(\xtpo_k).
$$
Notably, $\tilde s_k$ is a $\mathcal{T}_{\xtpo_k}$-valued Riemannian gradient.

The second change is to make the proposal on \Cref{alg:TDS} \Cref{alg_line:proposal} a tangent normal, as defined in \cref{eqn:likelihood_approx_riemannian_supp}.